\section{Introducción}

\subsection{Contexto y problema industrial}

Los robots móviles autónomos (AMR) se han consolidado como infraestructura intralogística para trasladar materiales entre estaciones, almacenes y líneas de producción. En 2024 se vendieron $102\,900$ robots de servicio profesional para transporte y logística, más de la mitad del total registrado en esa categoría por la International Federation of Robotics \citep{ifr2025service}. El despliegue a gran escala suele apoyarse en cargas, interfaces y recorridos normalizados. El sistema Kiva, por ejemplo, coordina cientos de vehículos mediante estanterías móviles, estaciones fijas y gestión central del tráfico \citep{wurman2008coordinating}.

Esa regularidad deja de estar disponible cuando una instalación debe mover bastidores voluminosos, piezas frágiles, subconjuntos pesados o cargas cuyo centro de masa exige varios apoyos. Las taxonomías de asignación multi-robot muestran que la dificultad aumenta al combinar capacidades heterogéneas, tareas que requieren varios agentes y restricciones temporales \citep{gerkey2004formal,korsah2013taxonomy}. Dimensionar cada robot para la carga máxima reduce modularidad; excluir las cargas pesadas interrumpe el flujo automatizado. El transporte cooperativo ofrece una tercera opción: formar una coalición temporal cuya capacidad conjunta se adapte a cada carga \citep{tuci2018cooperative,farivarnejad2022multirobot,an2023cooperative}.

La coalición debe cumplir más que una cuota numérica. Los robots seleccionados necesitan cubrir la demanda, ocupar contactos compatibles, producir la fuerza y el torque requeridos, conservar margen para acelerar y frenar, evitar obstáculos y sustituir miembros cuando cambian la batería o la comunicación. Por ello, la unidad científica de este trabajo no es la asignación aislada, sino la transición desde una preferencia distribuida hasta una ejecución físicamente válida.




\subsection{Brecha de investigación}

La asignación multi-robot (MRTA) y las subastas distribuidas ofrecen mecanismos sólidos para decidir qué agente atiende cada tarea. CBBA, por ejemplo, coordina paquetes de tareas mediante pujas y consenso \citep{choi2009consensus}. Los métodos de formación de coaliciones representan tareas que exigen varios agentes, aunque la búsqueda combinatoria crece con rapidez y suele abstraer la geometría de contacto \citep{sandholm1999coalition,rahwan2015coalition}. En ambos casos, una decisión válida en el espacio de tareas puede aceptar una coalición que no genere el \emph{wrench} de la carga.

La literatura de control cooperativo estudia formaciones, seguimiento y manipulación de objetos, pero con frecuencia parte de una coalición o una configuración de contacto ya definida \citep{tuci2018cooperative,farivarnejad2022multirobot}. Entre ambos bloques queda una brecha operacional: MRTA decide quién participa y el controlador decide cómo moverse, sin una prueba común que determine si esa selección puede ejecutar la misión completa. El problema se agrava bajo información local, porque el cierre entero, la factibilidad mecánica y la recuperación pueden requerir datos que no están disponibles en un solo vecindario.

Las dinámicas poblacionales aportan una representación analítica para repartir preferencia entre cargas mediante diferencias de utilidad \citep{sandholm2010population,quijano2017population}. En control distribuido se han usado para optimización, asignación de recursos y coordinación sobre grafos \citep{barreiro2017distributed,martinezpiazuelo2022tcss}. Su ventaja es la trazabilidad del equilibrio y del potencial; su límite es igualmente importante: una distribución continua no constituye por sí sola una coalición entera ni físicamente ejecutable. Esta tesis estudia precisamente la interfaz entre ambos niveles.






\subsection{Pregunta de investigación y enfoque}

La pregunta de investigación es: ¿en qué condiciones una arquitectura basada en decisiones locales y cierre entero puede formar y ejecutar coaliciones multi-AMR para cargas heterogéneas sin confundir cardinalidad, capacidad, factibilidad mecánica, seguridad y conectividad?

La memoria denomina \emph{coalición física} a una selección que supera una cadena explícita de comprobaciones: cuórum entero, capacidad efectiva, residual \emph{wrench}, autoridad dinámica, seguridad durante el movimiento y recuperación ante eventos. Esta definición convierte cada capa en una afirmación falsable. Una coalición aceptada en un escalón puede fallar en el siguiente; ese fallo no se descarta, sino que identifica qué abstracción resultó insuficiente.

La evaluación combina dos tipos de evidencia. A0--FULL ejecuta seis prefijos acumulativos sobre la misma planta Euler--Lagrange y constituye la prueba integrada de la tesis. SP1--SP8 son campañas modulares con plantas y escalas propias: aíslan reclutamiento homogéneo, capacidad heterogénea, contactos y \emph{wrench}, movimiento, carga extendida, recuperación, comunicación y escala. Sus resultados explican mecanismos concretos, pero no se suman como si procedieran de una única arquitectura de extremo a extremo.



\subsection{Contribuciones y organización de la memoria}

El trabajo aporta cuatro resultados delimitados por el simulador planar:
\begin{enumerate}[label=A\arabic*.,leftmargin=2.2em]
  \item una definición operativa de coalición física que separa aceptación por etapa, falso positivo físico y éxito final;
  \item una formulación que conecta déficit, capacidad efectiva, geometría de contacto y reparto de \emph{wrench}, junto con una aplicaci\'on derivada de ISS pr\'actica bajo supuestos expl\'icitos;
  \item la campaña A0--FULL, con mundos emparejados y extensiones decididas solo por precisión, para medir la interacción entre selección, mecánica, seguridad y recuperación;
  \item una auditoría modular SP1--SP8 que conserva resultados favorables, negativos, exploratorios y no estimables, y delimita qué parte del rendimiento procede del juego, del cierre o de una guardia global.
\end{enumerate}
Los objetivos y las hipótesis traducen estas contribuciones en criterios verificables. La metodología define la unidad experimental y el contrato inferencial. El marco teórico relaciona el problema con MRTA, juegos poblacionales, control distribuido y mecánica de la carga. Los resultados presentan primero la campaña integrada y después la evidencia modular necesaria para interpretarla. Las conclusiones responden a los objetivos y fijan los límites de validez.
